\section{变量保存方式}

最开始写这类 CFD 程序时，很容易把所有数组都直接放在一个结构体里，例如
\texttt{q1}、\texttt{q2}、\texttt{q3}、\texttt{f1}、\texttt{f2}、\texttt{f3}。
这种写法能运行，但读代码时经常需要回头查注释。02 版改成了按物理量分组保存。

\subsection{守恒量}

守恒量用 \texttt{Conserved} 保存：
\begin{lstlisting}
typedef struct {
    double *rho;
    double *rhou;
    double *rhoE;
} Conserved;
\end{lstlisting}

它对应
\[
\bm Q=(\rho,\rho u,\rho E)^{\mathrm T}.
\]
这样写之后，\texttt{rho}、\texttt{rhou}、\texttt{rhoE} 的含义比较直观，
比 \texttt{q1/q2/q3} 更容易检查。

\subsection{原始量}

原始量用 \texttt{Primitive} 保存：
\begin{lstlisting}
typedef struct {
    double *rho;
    double *u;
    double *p;
    double *E;
    double *a;
} Primitive;
\end{lstlisting}

这些量不是单独推进的变量，而是每一步由守恒量算出来：
\[
u=\frac{\rho u}{\rho},
\qquad
p=(\gamma-1)\left(\rho E-\frac12 \rho u^2\right),
\qquad
a=\sqrt{\gamma p/\rho}.
\]
其中声速 \(a\) 用来计算 CFL 时间步，\(\rho,u,p\) 也会用于人工粘性系数的计算。

\subsection{通量}

通量用 \texttt{Flux} 保存：
\begin{lstlisting}
typedef struct {
    double *mass;
    double *momentum;
    double *energy;
} Flux;
\end{lstlisting}

它对应 Euler 方程中的
\[
\bm F(\bm Q)=
\begin{bmatrix}
\rho u\\
\rho u^2+p\\
u(\rho E+p)
\end{bmatrix}.
\]
这样写以后，质量通量、动量通量、能量通量可以直接从变量名判断。

\subsection{不同时间层}

MacCormack 格式中一个时间步内会出现多个中间状态。代码中用
\texttt{TimeLayers} 统一保存：
\begin{lstlisting}
typedef struct {
    Conserved current;
    Conserved filtered;
    Conserved predicted;
    Conserved corrected;
} TimeLayers;
\end{lstlisting}

含义如下：

\begin{center}
\begin{tabular}{lll}
\toprule
代码变量 & 数学符号 & 含义 \\
\midrule
\texttt{current} & \(\bm Q^n\) & 当前时间层 \\
\texttt{filtered} & \(\bar{\bm Q}^{\,n}\) & 人工粘性滤波后 \\
\texttt{predicted} & \(\tilde{\bm Q}^{\,n+1}\) & 预测步结果 \\
\texttt{corrected} & \(\bm Q^{n+1}\) & 校正步结果 \\
\bottomrule
\end{tabular}
\end{center}

每完成一个时间步，代码把 \texttt{corrected} 复制回 \texttt{current}：
\begin{lstlisting}
conserved_copy(solver, &solver->state.current,
               &solver->state.corrected);
\end{lstlisting}
这样下一步仍然从 \texttt{current} 开始，时间层关系比较明确。

\subsection{求解器总结构}

所有参数和数组最后放在 \texttt{Solver} 中：
\begin{lstlisting}
typedef struct {
    Grid1D grid;
    GasModel gas;
    RiemannInitialCondition initial;
    TimeControl time;
    ArtificialViscosity av;
    MacCormackControl maccormack;

    TimeLayers state;
    FluxLayers flux;
    Primitive primitive;
    double *nu;
} Solver;
\end{lstlisting}

这里 \texttt{grid} 管网格，\texttt{gas} 管 \(\gamma\)，
\texttt{time} 管时间步，\texttt{av} 管人工粘性，
\texttt{maccormack} 管差分方向。这样虽然结构体多了一些，
但每一块变量的用途比较明确。

\subsection{内存分配}

每类数组都有对应的分配和释放函数，例如：
\begin{lstlisting}
static int conserved_alloc(Conserved *q, size_t n);
static void conserved_free(Conserved *q);
\end{lstlisting}

这样做的目的是减少漏分配、漏释放数组的可能，也少写重复代码。
比如一个 \texttt{Conserved} 里面总是有三条数组，
分配和释放时就应该一起处理。
